\subsection{Лекция 4 (10.10.26) МКO для неортогональных сеток}
В тесте \cvar{poisson2-fvmex} из файла \ename{poisson_fvmex_test.cpp}
реализовано решение двумерного уравнения Пуассона методом МКО с расширенным набором точек коллокаций.
Использованы с граничные условия первого рода (см. теорию в \secref{sec:FVM}).

Решения сохраняются в vtk-файлы \ename{poisson2_fvmex.vtk}.
Отталкиваясь от этой реализации необходимо:

\begin{enumerate}
\item Реализовать учёт неортогональности сетки при сборке СЛАУ градиентным методом (\secref{sec:nonortho_fvm})
      с помощью явной итерационной поправки (\secref{sec:ortho_correction_assembly})
      при фиксированном числе итераций.
      Следует рассмотреть два метода вычисления градиента в центрах ячеек (\secref{sec:gradu_in_cells})
      и три метода вычисления касательной производной (\secref{sec:fvm_duda}).
\item Проиллюстрировать численное решение, полученное на неструктурированной сетке.
      Сравнить решение без учёта поправки с одним из вариантов с поправкой.
\item Для каждого из шести запрограммированных вариантов настроек
      следует нарисовать график сеточной сходимости с использованием неструктурированной сетки.
      На каждый график нужно поместить базовый вариант без учёта поправки,
      а так же несколько вариантов с поправкой при разном количестве выполненных итераций.
\end{enumerate}

\paragraph{Рекомендации к программированию}
\begin{itemize}
\item
Для построения скошенной сетки воспользоваться скриптом \ename{tetragrid.py}
(см. \secref{sec:prog_grid}).
\item
Класс для решения задачи с учётом поправки \cvar{Problem2DSkew} следует отнаследовать
от имеющегося решателя \cvar{Problem2D}. При этом в конструктор
помимо сетки следует передавать настройку для метода вычисления градиента
и метода вычисления касательной.
\item
Процедуру решения с учётом поправки на скошенность следует
организовать согласно пошаговому алгоритму из \secref{sec:ortho_correction_assembly}.
Для того, чтобы избежать лишних вычислений следует разбить решение на этапы
инициализации и итерации.
\item
Для вычисления градиентов на гранях нужно воспользоваться
одной из двух реализаций интерфейса \cvar{IFvmFaceGradient}
из файла \ename{cfd/fvm/fvm_gradient.hpp}.
Таким образом у нового класса должно быть поле
\begin{cppcode}
std::shared_ptr<IFvmFaceGradient> grad_computer_;
\end{cppcode}
которое нужно определить в конструктора класса как (на примере метода наименьших квадратов):
\begin{cppcode}
grad_computer_ = std::make_shared<LeastSquaresFvmFaceGradient>(*grid, ecol_);
\end{cppcode}
Тогда вычисление массива градиентов в центрах ячеек по заданному полю \cvar{u_} примет вид
\begin{cppcode}
std::vector<Vector> gradu = grad_computer_->compute(u_);
\end{cppcode}
\item
Работа с расширенным набором точек коллокации в тестовой программе
осуществляется в объекте \cvar{Promblem2D::ecol_}, принадлежащего классу \cvar{FvmExtendedCollocations}.
Из него, в частности, берутся:
\begin{itemize}
\item \cvar{size()} -- количество точек коллокации,
\item \cvar{points} -- координаты точек коллокации,
\item \cvar{tab_face_colloc()} -- таблица связности \quo{грань -- точка коллокации}
\item и т.д. Полный список методов смотри в файле \ename{fvm/fvm_extended_collocations.hpp}.
\end{itemize}
Нумерация точек коллокаций устроена так, что первые $N$ (по количеству ячеек) индексов соответствуют внутренним точкам,
оставшиеся $N^\Gamma$ -- граничные точки коллокации.
Пользуясь этими методами, а также методом \cvar{IGrid::face_normal(iface)}
можно посчитать векторы $\vec{c'}$ и $\vec a$ по формулам из \secref{sec:fvm_duda}.
Дополнительно понадобятся функции вычисления модуля вектора \cvar{vector_abs} и
скалярного произведения \cvar{dot_product}.
Например, подсчёт массива векторов $\vec a$ по методу поворота может проводится по следующей процедуре
\begin{cppcode}
std::vector<Vector> a;
for (size_t iface=0; iface < grid_->n_faces(); ++iface){
    auto [icolloc, jcolloc] = ecol_.tab_face_colloc(iface);
    Vector cij = ecol_.points[jcolloc] - ecol_.points[icolloc];
    Vector c_prime = cij / vector_abs(cij);
    a.push_back(grid_->face_normal(iface) - c_prime);
}
\end{cppcode}
\end{itemize}
